% =============================================================================
% БЫСТРЫЕ КОМАНДЫ ДЛЯ КОНСПЕКТОВ
% =============================================================================
% Добавляйте свои команды в этот файл для ускорения набора

% =============================================================================
% МАТЕМАТИКА: МНОЖЕСТВА
% =============================================================================

\newcommand{\N}{\mathbb{N}}           % Натуральные числа
\newcommand{\Z}{\mathbb{Z}}           % Целые числа
\newcommand{\Q}{\mathbb{Q}}           % Рациональные числа
\newcommand{\R}{\mathbb{R}}           % Вещественные числа
\newcommand{\C}{\mathbb{C}}           % Комплексные числа
\newcommand{\F}{\mathbb{F}}           % Поле

\newcommand{\set}[1]{\left\{#1\right\}}           % {x}
\newcommand{\setmid}[2]{\left\{#1\mid#2\right\}}  % {x | условие}
\newcommand{\card}[1]{\left|#1\right|}            % |A| мощность

% =============================================================================
% МАТЕМАТИКА: ОПЕРАЦИИ
% =============================================================================

\newcommand{\abs}[1]{\left|#1\right|}             % |x|
\newcommand{\norm}[1]{\left\|#1\right\|}          % ||x||
\newcommand{\floor}[1]{\left\lfloor#1\right\rfloor}
\newcommand{\ceil}[1]{\left\lceil#1\right\rceil}
\newcommand{\inner}[2]{\langle#1,#2\rangle}       % <x,y> скалярное произведение

\DeclareMathOperator{\sgn}{sgn}
\DeclareMathOperator{\Ker}{Ker}
\DeclareMathOperator{\Img}{Im}
\DeclareMathOperator{\rk}{rk}
\DeclareMathOperator{\tr}{tr}
\DeclareMathOperator{\diag}{diag}
\DeclareMathOperator{\const}{const}

% =============================================================================
% МАТЕМАТИКА: ПРЕДЕЛЫ И СУММЫ
% =============================================================================

\newcommand{\limn}{\lim_{n\to\infty}}             % lim n->inf
\newcommand{\limx}[1]{\lim_{x\to#1}}              % lim x->a
\newcommand{\sumn}{\sum_{n=1}^{\infty}}           % sum n=1 to inf
\newcommand{\sumk}{\sum_{k=1}^{n}}                % sum k=1 to n
\newcommand{\prodn}{\prod_{n=1}^{\infty}}         % prod n=1 to inf

% =============================================================================
% МАТЕМАТИКА: ПРОИЗВОДНЫЕ И ИНТЕГРАЛЫ
% =============================================================================

\newcommand{\dd}{\,\mathrm{d}}                    % дифференциал
\newcommand{\dv}[2]{\frac{\mathrm{d}#1}{\mathrm{d}#2}}    % d/dx
\newcommand{\ddv}[2]{\frac{\mathrm{d}^2#1}{\mathrm{d}#2^2}} % d²/dx²
\newcommand{\pdv}[2]{\frac{\partial#1}{\partial#2}}        % частная производная
\newcommand{\pddv}[2]{\frac{\partial^2#1}{\partial#2^2}}   % вторая частная

% =============================================================================
% МАТЕМАТИКА: СТРЕЛКИ И ОТНОШЕНИЯ
% =============================================================================

\newcommand{\then}{\Rightarrow}                   % =>
\newcommand{\onlyif}{\Leftarrow}                  % <=
\newcommand{\iff}{\Leftrightarrow}                % <=>
\newcommand{\xto}[1]{\xrightarrow{#1}}            % стрелка с надписью
\newcommand{\mto}{\mapsto}                        % |->

\newcommand{\eq}{\equiv}                          % тождественно равно
\newcommand{\neq}{\not\equiv}
\newcommand{\iso}{\cong}                          % изоморфизм

% =============================================================================
% МАТЕМАТИКА: СОКРАЩЕНИЯ
% =============================================================================

\newcommand{\eps}{\varepsilon}
\newcommand{\vphi}{\varphi}
\newcommand{\vthe}{\vartheta}
\newcommand{\es}{\varnothing}                     % пустое множество

\newcommand{\ol}[1]{\overline{#1}}                % черта сверху
\newcommand{\ul}[1]{\underline{#1}}               % черта снизу
\newcommand{\wt}[1]{\widetilde{#1}}               % тильда
\newcommand{\wh}[1]{\widehat{#1}}                 % шляпка

% =============================================================================
% ФОРМАТИРОВАНИЕ: ТЕКСТ
% =============================================================================

\newcommand{\term}[1]{\textbf{#1}}                % Термин (выделение)
\newcommand{\eng}[1]{\textenglish{#1}}            % Английский текст
\newcommand{\todo}[1]{{\color{red}\textbf{TODO:} #1}}  % Пометка TODO

% =============================================================================
% ФОРМАТИРОВАНИЕ: БЛОКИ
% =============================================================================

% Важная информация
\newcommand{\important}[1]{%
    \begin{thmbox}
    \textbf{Важно:} #1
    \end{thmbox}
}

% Примечание на полях
\newcommand{\sidenote}[1]{%
    \marginpar{\footnotesize\color{secondary}#1}
}

% Доказательство (короткая версия)
\newcommand{\pf}[1]{%
    \begin{proof}
    #1
    \end{proof}
}

% =============================================================================
% БЫСТРЫЕ ОКРУЖЕНИЯ
% =============================================================================

% Теорема в боксе
\newenvironment{thm}[1][]{%
    \begin{thmbox}
    \begin{theorem}[#1]
}{%
    \end{theorem}
    \end{thmbox}
}

% Определение в боксе
\newenvironment{defn}[1][]{%
    \begin{defbox}
    \begin{definition}[#1]
}{%
    \end{definition}
    \end{defbox}
}

% Лемма в боксе
\newenvironment{lem}[1][]{%
    \begin{thmbox}
    \begin{lemma}[#1]
}{%
    \end{lemma}
    \end{thmbox}
}

% =============================================================================
% СПИСКИ БЫСТРЫЕ
% =============================================================================

% Компактный itemize
\newenvironment{items}{%
    \begin{itemize}[noitemsep,topsep=0pt]
}{%
    \end{itemize}
}

% Компактный enumerate
\newenvironment{enum}{%
    \begin{enumerate}[noitemsep,topsep=0pt]
}{%
    \end{enumerate}
}

% =============================================================================
% ДОПОЛНИТЕЛЬНЫЕ КОМАНДЫ
% =============================================================================
% Добавляйте свои команды ниже:


